\documentclass[11pt]{article}

\usepackage[margin=1in]{geometry}
\usepackage{amsmath, amssymb}
\usepackage{xcolor}
\usepackage[skins,breakable]{tcolorbox}
\usepackage{hyperref}

\hypersetup{
    colorlinks=true,
    linkcolor=blue!60!black,
    urlcolor=blue!60!black,
    citecolor=blue!60!black
}

\newcommand{\SymPy}{SymPy}

% Public artifact repository; \bugbundle links a bug's bundle folder into it.
% TODO: set the public artifact repository URL before release.
\newcommand{\artifactrepo}{https://github.com/TODO-org/TODO-repo}
\newcommand{\bugbundle}[1]{\href{\artifactrepo/tree/main/bug-report/#1}{\nolinkurl{bug-report/#1/}}}

\newtcolorbox{counterexamplebox}{
  enhanced,
  colback=yellow!5,
  colframe=orange!70!black,
  arc=2mm,
  boxrule=0.8pt,
  left=4mm,
  right=4mm,
  top=3mm,
  bottom=3mm,
  before=\par\medskip\noindent,
  after=\par\medskip,
  before upper=\raggedright
}

\title{SymPy Bug Card}
\author{}
\date{}

\begin{document}

\maketitle

\section*{Bug 1: solveset accepts \texorpdfstring{$\operatorname{asin}(x)=\pi$}{asin(x) = pi}}

\begin{counterexamplebox}
\textbf{Task.} Solve the principal inverse-sine equation over the complex
numbers.
\[
  \operatorname{asin}(x)=\pi,\qquad x\in\mathbb{C}.
\]

\medskip
\textbf{SymPy's answer.}
\[
  \operatorname{solveset}(\operatorname{asin}(x)=\pi,\ x,\ \mathbb{C})=\{0\}.
\]

\medskip
\textbf{Correct answer.}
\[
  \varnothing.
\]

\medskip
\textbf{Explanation.} SymPy evaluates \(\operatorname{asin}(0)=0\), so the returned candidate is not a solution; since \(\pi\) lies outside the principal range of \(\operatorname{asin}\), the equation is unsatisfiable.

\medskip
\textbf{Likely root cause.} In \nolinkurl{sympy/solvers/solveset.py},
\texttt{\_invert\_complex} appears to apply \(\sin\) as a global inverse of
principal \(\operatorname{asin}\), without carrying the required range condition, and the
finite candidate is not rejected by a residual check.  This is a new instance
of a known failure mode.

\medskip
\textbf{Artifact (GitHub).} \bugbundle{bug-001-solveset-accepts-asin-x-pi}
\end{counterexamplebox}

\end{document}
